\section{Objectives}
\label{sec:intro_objectives}

Guided by the two gaps identified in Section~\ref{sec:intro_problem}, we have set out the following objectives for this thesis. These are the goals we are working toward over the coming months, not a checklist of work already completed; the exact design of each component is expected to change as our implementation and early experiments give us feedback that the literature alone cannot.

\begin{enumerate}
    \item To investigate whether representing candidate vulnerabilities as a graph with explicit prerequisite relationships, rather than as an unordered list, can meaningfully improve how an autonomous agent chooses what to explore next -- and to design and formalize such a structure, which we are currently calling a \textbf{Vulnerability Dependency Graph (VDG)}, to the point where it can be implemented and tested.
    \item To explore ways of separating what an agent has \emph{confirmed} about a target from what it is merely \emph{hypothesizing}, so that the quality of discovery and the quality of planning can eventually be evaluated somewhat independently of each other.
    \item To settle on a multi-layer orchestration structure for the overall agent -- drawing on the structural patterns we observed recurring across the strongest-performing systems in our survey -- and to work out at what level of detail this structure needs to be specified before it is implementable.
    \item To examine whether letting an agent retain and reuse strategies across multiple missions offers a genuine advantage, or whether it mainly introduces risk (e.g., a strategy that worked against one version of a technology being harmful against another).
    \item To identify which existing, published benchmarks are suitable for evaluating this kind of system, and to begin drafting an evaluation plan around them rather than constructing new benchmarks of our own.
\end{enumerate}
